%========================= LIMITES ET PERSPECTIVES ==========================
\chapter{Limites et perspectives}

\section{Limites actuelles du prototype}

Par souci de transparence, les limites du livrable sont énoncées clairement.
Elles découlent de choix assumés, cohérents avec la nature \textbf{prototype}
du projet.

\begin{itemize}
  \item \textbf{Prototype local} : l'application n'est pas déployée en
        production ; elle s'exécute en environnement de développement.
  \item \textbf{Données majoritairement mockées / en mémoire} :
        utilisateurs, annonceurs, magasins, campagnes existantes, créatives et
        paramètres de compte sont réinitialisés à chaque redémarrage du
        backend. Seuls les \textbf{brouillons de campagne} sont réellement
        persistés.
  \item \textbf{PostgreSQL non branché en production} : les modèles et les
        migrations existent et sont prêts, mais la base n'est pas connectée
        par défaut.
  \item \textbf{Authentification non \og production \fg{}} : l'authentification
        est \emph{simulée} (aucun mot de passe vérifié, aucun jeton de session
        réel) ; le contrôle d'accès n'est appliqué que côté frontend.
  \item \textbf{DCO simulé} : aucun fichier de créative n'est réellement stocké
        (métadonnées uniquement), et les \emph{landing pages} sont entièrement
        \textbf{simulées} (pas de route publique réelle).
  \item \textbf{Reporting sur données de démonstration} : les indicateurs et
        graphiques sont générés côté frontend, non encore agrégés depuis une
        base réelle.
  \item \textbf{Absence de connexion à un réseau publicitaire réel} : aucune
        intégration DSP / réseau d'annonces (RTB, OpenRTB) n'est en place ;
        ce n'était pas l'objet du projet.
  \item \textbf{Pas de tests automatisés} : la validation reste manuelle et
        documentée.
  \item \textbf{URL de magasin non cliquables} : les URL des données de
        démonstration sont volontairement rendues non cliquables pour éviter
        d'ouvrir des pages inexistantes.
\end{itemize}

\section{Perspectives d'évolution}

Le prototype a été conçu pour \textbf{faciliter son évolution}. Les axes
d'amélioration suivants découlent directement de l'état actuel :

\begin{enumerate}[label=\textbf{\arabic*.}]
  \item \textbf{Intégration complète de PostgreSQL} : renseigner la connexion,
        appliquer les migrations SQL existantes et remplacer progressivement
        les services \emph{mockés} par de vraies requêtes.
  \item \textbf{Authentification et autorisation réelles} : mots de passe
        hachés, session ou jeton (JWT), contrôle d'accès vérifié côté API.
  \item \textbf{Persistance réelle des créatives DCO} et des variantes
        générées (stockage de fichiers).
  \item \textbf{Reporting sur données réelles} : agrégation de vraies
        statistiques depuis la base.
  \item \textbf{Dockerisation} de l'application (frontend, backend, base) pour
        un environnement reproductible.
  \item \textbf{Intégration continue (CI/CD)} : build, tests et déploiement
        automatisés.
  \item \textbf{Tests automatisés} : unitaires côté backend et de bout en bout
        côté frontend.
  \item \textbf{Déploiement en production} : hébergement du frontend statique
        et du backend ASGI, avec une stratégie de sauvegarde de la base.
  \item \textbf{Moteur DCO réel} et \textbf{intégration à un réseau
        publicitaire / DSP} pour la diffusion effective des campagnes.
  \item \textbf{Analytique avancée} : attribution des visites en magasin,
        indicateurs de fréquentation, tableaux de bord enrichis.
\end{enumerate}

Ces perspectives dessinent une trajectoire réaliste, du \textbf{prototype
validé} vers une \textbf{version de production}, en tirant parti de
l'architecture déjà mise en place.
